\chapter{The Road Ahead}
\label{ch:roadahead}

With the frontend served from the cluster (Chapter~26), every
deployable the project has is deployed and the application is usable
in a browser. The remaining milestones each swap something hand-rolled
for the industrial mechanism --- and each makes most sense now that you
know what it replaces.

\section{Observability}

Beyond \code{kubectl logs} lies the standard trio: \textbf{metrics}
(Prometheus scraping the \code{/actuator/prometheus} endpoints every
service already exposes; Grafana dashboards over the RED trio ---
rate, errors, duration), \textbf{traces} (the Zipkin wiring already in
every config, currently sampled at zero in the cluster --- one env var
turns it on once a Zipkin pod exists), and \textbf{alerts}
(Alertmanager: a page when a service's error rate or a PVC's fullness
crosses a line). This lands via the \code{kube-prometheus-stack} Helm
chart --- the promised moment where you \emph{consume} Helm without
authoring charts (Chapter~20's decision box, honored). Strimzi can
export broker metrics into the same Prometheus with a few lines of
\code{metricsConfig} on the Kafka resource --- consumer-group lag is
the first Kafka alert worth having.

\begin{hardware}
The full stack costs $\sim$1.5\,GB and nontrivial CPU --- the budget
says it fits, but it is the first thing to shed under load. Droppable
observability is a homelab luxury; production teams size for it first.
\end{hardware}

\section{GitOps: Argo CD}

The pipeline \emph{pushes} deployments (Jenkins holds a kubeconfig and
runs kubectl). GitOps inverts the arrow: an in-cluster agent
(Argo CD) \emph{pulls} --- continuously comparing a git repository's
manifests with the cluster and converging. What that buys, in the
vocabulary you now have: the cluster credential never leaves the
cluster (Chapter~22's RBAC concern dissolves --- CI only pushes images
and commits a tag bump); deletions finally propagate (Chapter~20's
prune pitfall, solved by design); drift becomes visible (a hand-edited
Deployment shows as ``OutOfSync''); and rollback is \code{git revert}.
Jenkins keeps CI --- build, test, push, update the tag in git ---
Argo CD takes CD. The Jenkinsfile's deploy stage, this book's
Chapter~23, becomes a commit.

\section{Retiring the classic stack}

The deliberate redundancy of Chapters~2 and 3 can now be paid off,
service by service:

\begin{itemize}
  \item \textbf{Eureka out}: set the gateway's routes from
  \code{lb://course-service} to \code{http://course-service:8082} ---
  Service DNS and readiness-gated endpoints do everything the registry
  did here. Delete a deployable, free 512\,Mi.
  \item \textbf{Config-server out}: mount the \code{configurations/}
  files as a ConfigMap and point
  \code{spring.config.import} at the mounted directory
  (\code{optional:file:/config/}). Same files, one fewer service, and
  config changes become manifest changes flowing through GitOps.
\end{itemize}

Both migrations are résumé material precisely because they run
\emph{toward} the platform: knowing why Eureka exists \emph{and} why
Kubernetes obsoletes it is the difference between using either and
understanding both.

\section{Jenkins into the cluster}

Chapter~22's homelab inversion --- builds on the controller, Jenkins
outside k3s --- reverses cleanly once the pipeline is stable: the
Kubernetes plugin spawns an agent pod per build, and Jenkins itself
becomes a StatefulSet with \code{JENKINS\_HOME} on a PVC. The
migration touches every concept in Part~V at once, which is why it is
scheduled last.

\section{OpenShift, briefly}

The manifests were kept portable on purpose: Jib's non-root images
(Chapter~13) satisfy OpenShift's arbitrary-UID security model, so
deploying to a free OpenShift sandbox should need only the
Ingress-to-Route translation (Chapter~18). Running that experiment
once validates the portability claims this book has been making ---
and OpenShift vocabulary (Routes, SCCs, BuildConfigs) maps one-to-one
onto chapters you have read.

\section{Where this leaves you}

Strip the project names away and the inventory reads like a platform
job description: the property-resolution machinery that makes one
artifact run anywhere; logs, consumer groups and delivery semantics;
images, layers and registries; reconciliation, probes, RBAC, storage
classes and operators; pipelines, RBAC-scoped deploy identities, and
the GitOps inversion. The homelab constraints were not limitations ---
they were the syllabus: every ``on this hardware'' box forced a
decision that money usually hides.

\begin{quiz}
\begin{enumerate}
  \item Which metric would you alert on first for the Kafka path, and
  what does it measure in Chapter~10's terms?
  \item Push-based vs.\ pull-based deployment: locate the cluster
  credential in each, and explain which chapter's pitfall GitOps
  pruning resolves.
  \item Sketch the Eureka retirement: what changes in the gateway's
  config, which two Kubernetes mechanisms replace registration and
  liveness, and what can be deleted?
  \item Why do Jib's images make the OpenShift experiment cheap, and
  which object translation remains?
\end{enumerate}
\end{quiz}
